% =====================================================================
\chapitre{Contexte général du projet}
% =====================================================================

Ce premier chapitre situe le projet. Il présente le cadre dans lequel il a été mené et
l'organisation de notre binôme, puis le besoin auquel répond \equipay et la problématique
qui en découle. Il fixe enfin les objectifs, le périmètre retenu et la démarche de travail
adoptée.

\subsection{Cadre du projet}

Le tableau~\ref{tab:cadre} résume le cadre dans lequel ce projet a été mené.

\begin{table}[H]
\centering
\small
\begin{tabularx}{\textwidth}{|L{3.2cm}|Y|}
\hline
\rowcolor{entete}\textbf{Critère} & \textbf{Description} \\ \hline
Établissement & École Marocaine des Sciences de l'Ingénieur (EMSI), 5\textsuperscript{ème} année du
cycle ingénieur, filière IASDR, option Intelligence Artificielle. \\ \hline
Module & Génération Textuelle : grands modèles de langage, affinage, évaluation de textes générés
et mise en production des modèles. \\ \hline
Équipe & Binôme : Alae Lahbichi et Haytham Benladoul. \\ \hline
Encadrement & Mme Hasnââ CHAABI. \\ \hline
Matériel & Un MacBook Pro M3 Pro (18 Go de mémoire unifiée) pour le développement, l'inférence et
le suivi des expériences ; un PC équipé d'une carte NVIDIA RTX 5080 pour l'entraînement des
modèles ; un iPhone pour les tests en conditions réelles. \\ \hline
Contrainte principale & Aucun service payant : services gratuits ou exécutés en local, modèles
ouverts. \\ \hline
\end{tabularx}
\caption{Cadre du projet}
\label{tab:cadre}
\end{table}

\subsection{Organisation du binôme}

Nous avons mené le projet de bout en bout, de l'idée de départ à l'application installée sur
téléphone. Le travail a été découpé par module (application mobile, base de données,
backend, puis chacun des modèles) afin que chacun puisse avancer sur sa partie sans bloquer
l'autre. Les choix d'architecture, les jeux de données et les résultats d'évaluation ont en
revanche toujours été discutés à deux, et chaque modification importante a été relue par
l'autre membre du binôme avant d'être intégrée. Les points réguliers avec notre encadrante
ont servi à valider les grandes orientations, en particulier la méthode d'évaluation des
modèles.

\subsection{Le partage des dépenses : un besoin du quotidien}

Les dépenses partagées font partie de la vie de groupe : un week-end, un voyage d'études,
une colocation ou un simple dîner. Le calcul n'est pas difficile en soi, mais il devient
vite pénible lorsque les dépenses s'accumulent, que les parts ne sont pas égales ou que
certains ont payé dans une autre devise. Pour un voyage, la question ne s'arrête d'ailleurs
pas à l'argent : il faut aussi décider de ce que l'on visite, dans quel ordre et à quel
moment, ce qui se fait le plus souvent dans une conversation de groupe où les propositions
se perdent.

\subsection{Problématique}

Deux difficultés se dégagent. La première tient à la saisie : chaque dépense doit être
entrée à la main (montant, payeur, participants, répartition), si bien que beaucoup
d'utilisateurs abandonnent en cours de route et reviennent au calcul approximatif. La
seconde tient à l'organisation : rien ne relie les dépenses du voyage à son programme, et
aucun outil grand public n'aide un groupe à choisir ensemble les lieux à visiter puis à les
enchaîner sans perdre de temps.

Les modèles de langage récents permettent de répondre aux deux. Un modèle vision-langage
peut lire un ticket de caisse, un modèle de langage affiné peut comprendre une phrase libre
et en tirer une dépense structurée, et un modèle généraliste peut rédiger un programme
clair à partir d'un itinéraire calculé. La question est alors de savoir comment intégrer
ces modèles dans une application fiable, gratuite et utilisable au quotidien, tout en
gardant la maîtrise de leur qualité : c'est la problématique à laquelle répond ce projet.

\subsection{Objectifs visés et délimitation du périmètre}

L'objectif est de concevoir une application mobile de partage des dépenses qui réduit la
saisie au minimum grâce à l'IA, aide le groupe à organiser ses sorties, et s'appuie sur une
démarche MLOps complète pour l'entraînement et le suivi des modèles.

Le périmètre retenu est présenté dans la figure~\ref{fig:perimetre}. Il comprend
l'application mobile (iOS en priorité, Android prêt), la base de données et
l'authentification, le backend d'IA, l'entraînement des modèles et leur suivi dans MLflow.
Le paiement intégré, la publication sur les magasins d'applications, l'hébergement du
backend dans le cloud et les notifications sont laissés comme évolutions futures.

\begin{figure}[H]
\centering
\begin{tikzpicture}[
  titre/.style={rounded corners=4pt, minimum width=6.6cm, minimum height=0.9cm, font=\bfseries\color{white}},
  item/.style={rounded corners=3pt, minimum width=6.6cm, minimum height=0.72cm, font=\footnotesize, draw=#1!60, fill=#1!8},
  ]
  \node[titre, fill=lagon] (in) at (0,0) {Inclus dans le projet};
  \node[titre, fill=gristexte] (out) at (7.3,0) {Hors périmètre (évolutions)};
  \foreach \t [count=\i] in {{Application mobile Flutter (iOS, Android prêt)},
                             {Comptes, activités, dépenses et soldes},
                             {Base Supabase : RLS et temps réel},
                             {Assistant conversationnel},
                             {Scanner de tickets et saisie par phrase},
                             {Programme de voyage et trajet optimisé},
                             {Entraînement et suivi dans MLflow}}
    \node[item=lagon] at (0,-0.35-0.85*\i) {\t};
  \foreach \t [count=\i] in {{Paiement intégré entre membres},
                             {Publication App Store et Play Store},
                             {Backend d'IA hébergé dans le cloud},
                             {Notifications push},
                             {Connexion avec un compte Apple}}
    \node[item=gristexte] at (7.3,-0.35-0.85*\i) {\t};
\end{tikzpicture}
\caption{Périmètre fonctionnel retenu : fonctionnalités incluses et exclues}
\label{fig:perimetre}
\end{figure}

\subsection{Démarche de travail et organisation du projet}

Nous avons adopté une démarche itérative et incrémentale, inspirée des méthodes agiles :
chaque cycle court se termine par une fonctionnalité testable sur téléphone. Comme le
montre la figure~\ref{fig:phases}, le projet s'est déroulé en cinq phases. Le cadrage a
fixé les besoins, le périmètre et les choix techniques. La conception a porté sur le schéma
de la base, l'architecture et la charte graphique. Le développement a produit l'application
et le backend. L'entraînement des modèles s'est déroulé en parallèle : préparation des
données, affinage sur GPU, évaluation et conversion. Les tests et la finalisation ont
regroupé la validation sur iPhone, l'optimisation et la rédaction de ce rapport.

\begin{figure}[H]
\centering
\begin{tikzpicture}[
  phase/.style={circle, minimum size=1.95cm, draw=white, line width=3pt, fill=#1, align=center,
                font=\footnotesize\bfseries, text=white, drop shadow={opacity=0.18}},
  bulle/.style={rectangle, rounded corners=6pt, draw=black!15, fill=white, align=center,
                font=\scriptsize, text width=2.75cm, minimum height=1.1cm},
  ]
  \fill[grisfond, rounded corners=6pt] (-1.6,-3.6) rectangle (14.1,3.6);
  \foreach \x/\y/\c/\n/\t/\d [count=\i] in {
      0/0.75/lagon/01/Cadrage/{Besoins, objectifs,\\ choix techniques},
      3.1/-0.75/violet/02/Conception/{Base de données,\\ architecture,\\ charte graphique},
      6.25/0.75/orange/03/Développement/{Application mobile,\\ backend, intégration\\ des modèles},
      9.4/-0.75/bleuclair/04/{Entraînement\\ des modèles}/{Données, affinage\\ LoRA, évaluation,\\ conversion MLX},
      12.5/0.75/recevoir/05/{Tests et\\ finalisation}/{Validation sur iPhone,\\ optimisation,\\ rapport}}
  {
    \node[phase=\c] (p\i) at (\x,\y) {\n\\[1pt]\t};
    \ifodd\i
      \node[bulle] (b\i) at (\x,-2.55) {\d};
      \draw[\c, thick] (p\i.south) -- (b\i.north);
      \fill[\c] (p\i.south) circle (2pt);
    \else
      \node[bulle] (b\i) at (\x,2.55) {\d};
      \draw[\c, thick] (p\i.north) -- (b\i.south);
      \fill[\c] (p\i.north) circle (2pt);
    \fi
  }
\end{tikzpicture}
\caption{Les cinq phases de la démarche de travail adoptée}
\label{fig:phases}
\end{figure}

\subsection{Conclusion du chapitre}

Ce chapitre a présenté le cadre du projet, l'organisation de notre binôme et le besoin
auquel répond \equipay. La problématique posée, les objectifs, le périmètre retenu et la
démarche adoptée constituent le point de départ de l'étude préalable présentée dans le
chapitre suivant.
